% Esquisse d'un programme (cote 119, copy 1): archive cote 119, site page 10 (PDF page 11), typescript page 8.
% typescript. Draft compared: EsquisseFr.pdf (text layer):275-361.
% AI-transcribed by Claude Opus 5.5 (claude-opus-5-5, adaptive thinking, effort high), Message Batches API, batch esq1 (msgbatch_01LNsogFyF3y6KSp9UbtJtBk), 2026-11-12; not yet checked by eye.
\scanpage{119/10}
\mspage{8}

de l'étage O), on trouve \uline{douze} ``\uline{pièces de construction}'' \uline{fondamentales}
(6 de genre O, 6 de genre 1) dans un ``jeu de Légo-Teichmüller (grande
boite), où les points marqués sur les surfaces envisagées sont rem-%
placés par des ``trous'' à bord, de façon à avoir des surfaces à bord,
donc des pièces de construction qui peuvent s'assembler par frotte-%
ment doux comme dans le jeu de Légo ordinaire cher à nos enfants
(ou petits-enfants...). Par assemblage on trouve un moyen tout ce
qu'il y a de visuel pour construire tout type de surface (ce sont
ces assemblages essentiellement qui seront les ``points bases'' pour
notre fameuse tour), et aussi de visualiser les ``\uline{chemins}'' \uline{élémentaires}
par des opérations tout aussi concrètes tels des ``twists'', ou des
automorphismes des pièces du jeu, et d'écrire les \uline{relations fondamentales}
entre chemins composés. Suivant la taille (et le prix !) de la boite
de construction utilisée, on trouve d'ailleurs de nombreuses descrip-%
tions différentes de la tour de Teichmüller par générateurs et rela-%
tions. La boite la plus petite est réduite à des pièces toutes iden-%
tiques, de type (O,3) - ce sont les ``pantalons'' de Thurston, et le
jeu de Légo-Teichmüller que j'essaie de décrire, issu de motivations
et de réflexions de géométrie algébrique absolue sur le corps \hand{$\underline{Q}$},
est très proche du jeu de ``chirurgie géodésique hyperbolique'' de
Thurston, dont j'ai appris l'existence l'an dernier par Yves Ladegaillerie.
Dans un microséminaire avec Carlos Contou-Carrère et Yves Ladegraillerie,
nous avons amorcé une réflexion dont un des objets est de confronter
les deux points de vue, qui se complètent mutuellement.

J'ajoute que chacune des douze pièces de construction de
la ``grande boite'' se trouve muni d'une décomposition cellulaire cano-%
nique, stable par toutes les symétries, ayant comme seuls sommets les
``points marqués'' (ou centres des trous), et comme arètes certains che-%
mins géodésiques (pour la structure riemanienne canonique sur la sphère
ou le tore envisagé) entre certaines paires de sommets (savoir ceux
qui se trouvent sur un même ``lieu réel'', pour une structure réelle
convenable de la courbe algébrique complexe envisagée). Par suite,
toutes les surfaces obtenues dans ce jeu par assemblage sont munies
de structures cellulaires canoniques, qui à leur tour (cf.§3 plus bas)

\margin{10}
